%% kapa: exam=1 round=35 year=2024 period=1 form=A subject=civil_law subject_ko=민법 seq=38 number=38
%% source: Q-Net 감정평가사 기출문제 (cst00309) artlSeq=5213229 file=raw/1008/5213229/1교시 문제지 원본.pdf
%% conversion: kordoc 4.24.51 -> tools/gen1.py (auto); review=proofread-sonnet; flags=none
\begin{question}{38}
저당권의 효력이 미치는 범위에 관한 설명으로 옳지 않은 것은? (다툼이 있으면 판례에 따름)
\begin{choices}
  \item 담보권 실행을 위하여 저당부동산을 압류한 경우, 저당부동산의 압류 이후 발생한 차임채권에는 저당권의 효력이 미친다.
  \item 주물 그 자체의 효용과는 직접 관계없지만 주물 소유자의 상용에 공여되고 있는 물건이 경매목적물로 평가되었다면 경매의 매수인이 소유권을 취득한다.
  \item 구분건물의 전유부분에 대한 저당권의 효력은 특별한 사정이 없는 한 대지사용권에도 미친다.
  \item 기존건물에 부합된 증축부분이 기존건물에 대한 경매절차에서 경매목적물로 평가되지 아니하였더라도 경매의 매수인이 증축부분의 소유권을 취득한다.
  \item 특약이 없는 한 건물에 대한 저당권의 효력은 건물의 소유를 목적으로 하는 지상권에도 미친다.
\end{choices}
\end{question}
\answer{2}
